\section{Constant-Term Normalisation for Primitive CHL Borcherds Denominators}
\label{sec:cn-zero-primitive-chl-normalisation}

\providecommand{\cO}{\mathcal{O}}
\providecommand{\kBKM}{\kappa_{\mathrm{BKM}}}
\providecommand{\kch}{\kappa_{\mathrm{ch}}}
\providecommand{\kcat}{\kappa_{\mathrm{cat}}}
\providecommand{\kfib}{\kappa_{\mathrm{fiber}}}

\paragraph{Convention.}
Let
\[
  \varphi_N(\tau,z)
  =\sum_{n,\ell} c_N(n,\ell)\,q^n\zeta^\ell,
  \qquad q=e^{2\pi i\tau},\quad \zeta=e^{2\pi iz},
\]
be the weight-$0$, index-$1$ Jacobi input for the primitive
CHL/BKM Borcherds denominator at order
$N\in\{1,2,3,4,6\}$.  The constant
\[
  c_N(0):=c_N(0,0)=[q^0\zeta^0]\,\varphi_N(\tau,z)
\]
is the coefficient entering Borcherds' singular-theta weight theorem.
Thus
\[
  \kBKM(\Phi_N)=\mathrm{wt}(\Phi_N)=\frac{c_N(0)}{2}.
\]
This is the automorphic-weight invariant of the Borcherds product,
distinct from $\kch$, $\kcat=\chi(\cO_X)$, and $\kfib$.

\begin{theorem}[Primitive CHL constant terms]
\label{thm:cn-zero-primitive-chl-normalisation}
For the primitive CHL/BKM denominator ladder
$N\in\{1,2,3,4,6\}$,
\[
  \begin{array}{c|ccccc}
  N & 1 & 2 & 3 & 4 & 6 \\ \hline
  c_N(0) & 10 & 8 & 6 & 4 & 2 \\
  \kBKM(\Phi_N) & 5 & 4 & 3 & 2 & 1
  \end{array}
\]
\end{theorem}

\begin{proof}
The Borcherds weight theorem gives
\[
  \mathrm{wt}\,\operatorname{Bor}(\varphi_N)
  =\frac{[q^0\zeta^0]\varphi_N}{2}.
\]
Therefore the constant-term table is twice the paramodular weight table
once the primitive denominator form has been fixed.

At $N=1$, the Eichler-Zagier generator satisfies
\[
  \phi_{0,1}^{\mathrm{EZ}}(\tau,z)
  =\zeta^{-1}+10+\zeta+O(q).
\]
The monic Gritsenko-Nikulin denominator is
\[
  D_5:=64^{-1}\Delta_5,\qquad
  [q^{1/2}r^{1/2}s^{1/2}]\Delta_5=64.
\]
The primitive Borcherds input is \(\phi_{0,1}^{\mathrm{EZ}}\), so
\(c_1(0)=10\) and
\[
  \operatorname{Bor}(\phi_{0,1}^{\mathrm{EZ}})=D_5,\qquad
  \mathrm{wt}(\Delta_5)=10/2=5.
\]
The K3 elliptic genus is
\[
  Z_{K3}=2\phi_{0,1}^{\mathrm{EZ}}.
\]
The doubled input lies on the scalar-square line:
\[
  \operatorname{Bor}(2\phi_{0,1}^{\mathrm{EZ}})=D_5^2,\qquad
  \Phi_{10}^{\theta}=\Delta_5^2,\qquad
  \Phi_{10}^{\mathrm{OP}}=D_5^2=4096^{-1}\Delta_5^2.
\]
The primitive denominator line is \(D_5=64^{-1}\Delta_5\), with
automorphic weight \(5\).

For \(N\in\{2,3,4,6\}\), Gritsenko-Nikulin's rank-three
paramodular construction gives primitive denominator weights
\[
  k(2)=4,\qquad k(3)=3,\qquad k(4)=2,\qquad k(6)=1.
\]
Borcherds' formula gives \(c_N(0)=2k(N)\) for these four orders,
recovering the four nontrivial entries in the theorem table.

The same constants are recovered from the theta-component refinement of
the twined K3 elliptic genus.  Let
\[
  T_N=\mathrm{tr}\!\left(g_N^\ast\mid H^\ast(K3,\mathbb Q)\right)
\]
for a symplectic K3 automorphism \(g_N\), and let \(A_N\) be the
odd theta-component constant in the index-$1$ theta decomposition of
\(\varphi_N\).  The \(q^0\)-specialisation gives
\[
  T_N=c_N(0)+2A_N.
\]
For the five CHL orders,
\[
\begin{array}{c|ccccc}
N & 1 & 2 & 3 & 4 & 6 \\ \hline
T_N & 24 & 16 & 12 & 10 & 8 \\
A_N & 7 & 4 & 3 & 3 & 3 \\
c_N(0)=T_N-2A_N & 10 & 8 & 6 & 4 & 2 \\
\kBKM(\Phi_N)=c_N(0)/2 & 5 & 4 & 3 & 2 & 1
\end{array}
\]
which agrees with the paramodular weight computation.
\end{proof}

\begin{remark}[Coefficient conventions separated]
\label{rem:cn-zero-coefficient-conventions-separated}
The singly-twined \(M_{23}\)-style coefficient column has
\[
\begin{array}{c|ccccc}
N & 1 & 2 & 3 & 4 & 6 \\ \hline
c_N^{\mathrm{tw}}(0) & 10 & 4 & 2 & 2 & 2 \\
c_N^{\mathrm{tw}}(0)/2 & 5 & 2 & 1 & 1 & 1
\end{array}
\]
The primitive CHL denominator weights are the five weights displayed in
Theorem~\ref{thm:cn-zero-primitive-chl-normalisation}.  Borcherds'
weight formula separates these two coefficient columns before any
geometric interpretation enters.

The Gritsenko-Cl\'ery diagonal-divisor catalogue has eight rows:
\[
\begin{array}{c|cccccccc}
(t,N;k) &
(1,1;5)&(2,1;2)&(1,2;3)&(3,1;1)&
(1,3;2)&(4,1;\tfrac12)&(1,4;\tfrac32)&(2,2;1)\\ \hline
2k & 10&4&6&2&4&1&3&2 .
\end{array}
\]
The indexing set is the paramodular triple \((t,N;k)\).  The primitive
CHL denominator uses the order set \(N\in\{1,2,3,4,6\}\).  The two
systems meet at \((t,N)=(1,1)\), where both constructions give
\(\Delta_5\) and weight \(5\).

For the order-indexed Nikulin products at \(N\in\{5,7,8\}\), the cover
group, twined Jacobi input, and Borcherds normalisation are part of the
data.  With those inputs fixed, the constants are
\[
\begin{array}{c|ccc}
N & 5 & 7 & 8 \\ \hline
c_N(0) & 4 & 2 & 2 \\
\mathrm{wt}(\Phi_N)=c_N(0)/2 & 2 & 1 & 1
\end{array}
\]
Borcherds products exist for these orders.  A smooth \(K3\times E\)
CHL Calabi-Yau threefold host requires an origin-fixing automorphism of
the elliptic factor, and the cyclotomic condition \(\varphi(N)\mid 2\)
selects \(N\in\{1,2,3,4,6\}\).
\end{remark}

\begin{remark}[Four invariant discipline]
\label{rem:cn-zero-four-invariant-discipline}
The identity
\[
  \kBKM(\Phi_N)=c_N(0)/2
\]
is independent of the additive expression
\[
  \kch(\mathcal A_{X_N})+\chi(\cO_{\mathrm{fiber}}).
\]
At \(N=1\), the Borcherds side is \(5\); the compact total-space Hodge
supertrace on \(K3\times E\) and the elliptic-fibre Euler characteristic
give \(0+0\).  At \(N=2\), the Borcherds side is \(4\); the corresponding
additive expression gives \(1\).  The four invariants \(\kch\), \(\kcat\),
\(\kBKM\), and \(\kfib\) answer different questions and have different
functorialities.
\end{remark}

\paragraph{References.}
Borcherds 1998, Theorem~13.3, supplies the weight formula
\(\mathrm{wt}\,\operatorname{Bor}(\varphi)=c(0,0)/2\).
Eichler-Zagier fixes
\(\phi_{0,1}^{\mathrm{EZ}}=\zeta^{-1}+10+\zeta+O(q)\).
Gritsenko-Nikulin's rank-three paramodular construction fixes the
primitive weights displayed in Theorem~\ref{thm:cn-zero-primitive-chl-normalisation}.
The theta-component table above is the same computation in the
twined K3 elliptic-genus normalisation.
